\section{Feynman 路径积分}
\subsection{时间切片与相空间路径积分}
把区间 $\left[t_0,t\right]$ 分为 $N$ 段，令
\begin{equation}
\varepsilon=\frac{t-t_0}{N},
\qquad t_j=t_0+j\varepsilon,
\qquad x_j=x\left(t_j\right),
\qquad x_N=x,
\qquad x_0=x_0.
\end{equation}
反复使用传播子的合成律，
\begin{equation}
K\left(x,t;x_0,t_0\right)=
\lim_{N\to\infty}\int\prod_{j=1}^{N-1}\mathrm{d}x_j
\prod_{j=0}^{N-1}\Braket{x_{j+1}|\mathrm{e}^{-\mathrm{i}\varepsilon\hat H/\hbar}|x_j}.
\label{eq:time-sliced-K}
\end{equation}
在每个短时间片插入动量完备性，并采用统一的离散化约定，可写成
\begin{equation}
\Braket{x_{j+1}|\mathrm{e}^{-\mathrm{i}\varepsilon\hat H/\hbar}|x_j}
=\int_{-\infty}^{\infty}\frac{\mathrm{d}p_j}{2\pi\hbar}
\exp\left\{\frac{\mathrm{i}\varepsilon}{\hbar}
\left[p_j\frac{x_{j+1}-x_j}{\varepsilon}-H\left(p_j,x_j,t_j\right)\right]\right\}
+O\left(\varepsilon^2\right).
\end{equation}
连续极限的形式记号为
\begin{equation}
\boxed{K\left(x,t;x_0,t_0\right)=
\int\mathcal D x\,\mathcal D p\,
\exp\left\{\frac{\mathrm{i}}{\hbar}
\int_{t_0}^{t}\left[p\dot x-H\left(p,x,t'\right)\right],\mathrm{d}t'\right\}.}
\label{eq:phase-space-path-integral}
\end{equation}
路径满足固定端点 $x\left(t_0\right)=x_0$、$x\left(t\right)=x$。式中的测度由有限维时间切片极限定义，不能把 $\mathcal D x$、$\mathcal D p$ 当作普通微分量。

\subsection{位形空间路径积分}
对
\begin{equation}
H\left(p,x,t\right)=\frac{p^2}{2m}+V\left(x,t\right),
\end{equation}
每个 $p_j$ 积分都是 Gaussian 积分。短时间传播子为
\begin{equation}
K\left(x_{j+1},t_{j+1};x_j,t_j\right)
=\sqrt{\frac{m}{2\pi\mathrm{i}\hbar\varepsilon}}
\exp\left\{\frac{\mathrm{i}\varepsilon}{\hbar}
\left[\frac{m}{2}\left(\frac{x_{j+1}-x_j}{\varepsilon}\right)^2
-V\left(x_j,t_j\right)\right]\right\}+O\left(\varepsilon^2\right).
\label{eq:short-time-kernel}
\end{equation}
代入式\eqref{eq:time-sliced-K}，得到严格的离散定义
\begin{align}
K\left(x,t;x_0,t_0\right)
&=\lim_{N\to\infty}
\left(\frac{m}{2\pi\mathrm{i}\hbar\varepsilon}\right)^{N/2}
\int_{-\infty}^{\infty}\prod_{j=1}^{N-1}\mathrm{d}x_j\notag\\
&\quad\times
\exp\left\{\frac{\mathrm{i}\varepsilon}{\hbar}\sum_{j=0}^{N-1}
\left[\frac{m}{2}\left(\frac{x_{j+1}-x_j}{\varepsilon}\right)^2
-V\left(x_j,t_j\right)\right]\right\}.
\label{eq:configuration-discrete}
\end{align}
引入 Lagrange 量与作用量
\begin{equation}
L\left(\dot x,x,t\right)=\frac{m}{2}\dot x^2-V\left(x,t\right),
\qquad
S\left[x\right]=\int_{t_0}^{t}L\left(\dot x,x,t'\right)\,\mathrm{d}t',
\end{equation}
便得到位形空间路径积分的简写
\begin{equation}
\boxed{K\left(x,t;x_0,t_0\right)=
\int_{x\left(t_0\right)=x_0}^{x\left(t\right)=x}\mathcal D x\,
\mathrm{e}^{\mathrm{i}S\left[x\right]/\hbar}.}
\label{eq:configuration-path-integral}
\end{equation}
它表示从初始点到终点的所有路径都贡献一个模为一、相位为 $S/\hbar$ 的振幅，传播子是这些振幅的相干叠加。

\subsection{由短时间核恢复 Schr\"odinger 方程}
令初始时刻为 $t$、终止时刻为 $t+\varepsilon$。由短时间核，
\begin{align}
\psi\left(x,t+\varepsilon\right)
&=\sqrt{\frac{m}{2\pi\mathrm{i}\hbar\varepsilon}}
\int_{-\infty}^{\infty}\mathrm{d}\xi\,
\exp\left(\frac{\mathrm{i}m\xi^2}{2\hbar\varepsilon}\right)
\exp\left[-\frac{\mathrm{i}\varepsilon}{\hbar}V\left(x+\xi,t\right)\right]
\psi\left(x+\xi,t\right)+O\left(\varepsilon^2\right),
\label{eq:short-step-wavefunction}
\end{align}
其中令 $x'=x+\xi$。快速振荡使主要贡献来自 $\xi=O\left(\sqrt\varepsilon\right)$。展开到所需阶数：
\begin{align}
\psi\left(x+\xi,t\right)
&=\psi+\xi\frac{\partial\psi}{\partial x}
+\frac{\xi^2}{2}\frac{\partial^2\psi}{\partial x^2}+O\left(\xi^3\right),\\
\exp\left[-\frac{\mathrm{i}\varepsilon}{\hbar}V\left(x+\xi,t\right)\right]
&=1-\frac{\mathrm{i}\varepsilon}{\hbar}V\left(x,t\right)+o\left(\varepsilon\right).
\end{align}
奇次 Gaussian 矩为零，而
\begin{align}
\int_{-\infty}^{\infty}\exp\left(\frac{\mathrm{i}m\xi^2}{2\hbar\varepsilon}\right)\,\mathrm{d}\xi
&=\sqrt{\frac{2\pi\mathrm{i}\hbar\varepsilon}{m}},\\
\int_{-\infty}^{\infty}\xi^2
\exp\left(\frac{\mathrm{i}m\xi^2}{2\hbar\varepsilon}\right)\,\mathrm{d}\xi
&=\frac{\mathrm{i}\hbar\varepsilon}{m}
\sqrt{\frac{2\pi\mathrm{i}\hbar\varepsilon}{m}}.
\end{align}
保留到 $O\left(\varepsilon\right)$，与
$\psi\left(x,t+\varepsilon\right)=\psi+\varepsilon\partial_t\psi+o\left(\varepsilon\right)$ 比较，得到
\begin{equation}
\mathrm{i}\hbar\frac{\partial\psi}{\partial t}
=\left[-\frac{\hbar^2}{2m}\frac{\partial^2}{\partial x^2}+V\left(x,t\right)\right]\psi.
\end{equation}
因此正确归一化的 Feynman 短时间核与 Schr\"odinger 方程等价。

\subsection{经典极限与驻相路径}
当典型作用量远大于 $\hbar$ 时，相邻路径的相位通常快速相消。只有作用量在一阶变分下不变的路径附近能够相干叠加：
\begin{equation}
\delta S\left[x_{\mathrm c}\right]=0.
\end{equation}
令端点固定，即 $\delta x\left(t_0\right)=\delta x\left(t\right)=0$，分部积分得
\begin{align}
\delta S
&=\int_{t_0}^{t}\left[
\frac{\partial L}{\partial x}\delta x
+\frac{\partial L}{\partial\dot x}\delta\dot x\right],\mathrm{d}t'\notag\\
&=\int_{t_0}^{t}\left[
\frac{\partial L}{\partial x}-\frac{\mathrm{d}}{\mathrm{d}t'}
\left(\frac{\partial L}{\partial\dot x}\right)\right]\delta x\,\mathrm{d}t'.
\end{align}
由于 $\delta x$ 任意，经典路径满足 Euler--Lagrange 方程
\begin{equation}
\boxed{\frac{\mathrm{d}}{\mathrm{d}t}
\left(\frac{\partial L}{\partial\dot x}\right)-\frac{\partial L}{\partial x}=0.}
\end{equation}
这说明经典运动是半经典极限中的驻相路径。

\subsection{两个经典作用量}
\begin{example}[自由粒子]
固定端点 $\left(x_1,t_1\right)$、$\left(x_2,t_2\right)$，自由粒子的经典路径为
\begin{equation}
x_{\mathrm c}\left(t\right)=x_1+
\frac{x_2-x_1}{t_2-t_1}\left(t-t_1\right).
\end{equation}
代入 $L=m\dot x^2/2$，得到
\begin{equation}
\boxed{S_{\mathrm c}=
\frac{m\left(x_2-x_1\right)^2}{2\left(t_2-t_1\right)}.}
\end{equation}
式\eqref{eq:free-propagator}的指数正是 $\mathrm{i}S_{\mathrm c}/\hbar$。
\end{example}

\begin{example}[一维谐振子]
令 $T=t_2-t_1$。满足两个端点条件的经典路径可写成
\begin{equation}
x_{\mathrm c}\left(t\right)=
\frac{x_2\sin\left[\omega\left(t-t_1\right)\right]
+x_1\sin\left[\omega\left(t_2-t\right)\right]}
{\sin\left(\omega T\right)}.
\end{equation}
利用运动方程 $\ddot x_{\mathrm c}+\omega^2x_{\mathrm c}=0$，经典作用量化为边界项：
\begin{align}
S_{\mathrm c}
&=\frac{m}{2}\int_{t_1}^{t_2}
\left(\dot x_{\mathrm c}^{\,2}-\omega^2x_{\mathrm c}^{\,2}\right)\,\mathrm{d}t
=\frac{m}{2}\left.x_{\mathrm c}\dot x_{\mathrm c}\right|_{t_1}^{t_2}\notag\\
&=\boxed{\frac{m\omega}{2\sin\left(\omega T\right)}
\left[\left(x_1^2+x_2^2\right)\cos\left(\omega T\right)-2x_1x_2\right].}
\end{align}
式\eqref{eq:oscillator-propagator}同样含有相位 $\mathrm{e}^{\mathrm{i}S_{\mathrm c}/\hbar}$；对二次作用量，绕经典路径的涨落积分仍是 Gaussian 积分，因而半经典形式在此恰好是精确结果。
\end{example}
